\section{Notation}
\label{app:notation}

\begin{table*}[t]
\centering\footnotesize
\setlength{\tabcolsep}{3pt}
\renewcommand{\arraystretch}{1.08}
\begin{tabular}{@{}l>{\raggedright\arraybackslash}p{4.2cm}>{\raggedright\arraybackslash}p{8.9cm}@{}}
\toprule
Symbol & Space / dimension & Meaning \\
\midrule
\multicolumn{3}{@{}l}{\emph{Input and Thinker (\S\ref{sec:method-task}--\ref{sec:method-simul})}} \\
$\bm{x}$, $T$ & $\R^{16000T}$; seconds & recording (16\,kHz samples) and its duration \\
$\bm{A}=[\bm{a}_1;\dots;\bm{a}_N]$ & $\R^{N\times d}$ & acoustic tokens: row $i$ = 80\,ms of audio centred at $\bar{t}_i$; $d=2048$ Thinker hidden dimensions \\
$\bm{P}$, $P_{f,k}$ & $[0,1]^{F\times S}$ & diarizer activities: row $f$ = 80\,ms frame at $t_f$, column $k$ = one of $S=8$ speaker slots \\
$\bm{s}_j$ & $\R^{d}$ & speaker token of the $j$-th block of $K=5$ acoustic tokens (0.4\,s) \\
$\bm{E}$ & $\R^{(N+M)\times d}$ & interleaved input, $M=\lceil N/K\rceil$ speaker tokens after their blocks $\bm{A}_j\in\R^{K\times d}$ \\
$\bm{\iota}$ & token sequence & text instruction \\
$\bm{y}$, $y_u$ & $\mathcal{V}^{|\bm{y}|}$ & serialised output; $\mathcal{V}$ = text vocabulary + 251 time tokens + 8 speaker tokens \\
$k_n$, $a_n,b_n$, $\bm{y}^{(n)}$ & $\{1,\dots,S\}$; $[0,T]$ (s); $\mathcal{V}^{|\bm{y}^{(n)}|}$ & segment $n$: speaker slot, source start and end, translation \\
$\alpha_n,\beta_n$ & $\{0,\dots,250\}$ & segment start and end in 80\,ms steps (time tokens $\tim{\alpha_n},\tim{\beta_n}$) \\
$\omega(y_u)$, $\Omega$ & $\{1,1.5,2\}$; $\R_{>0}$ & token weight (other / time / speaker) and its sum \\
$\Delta$, $\gamma$, $\epsilon$ & seconds & read interval (1), diarizer finalisation delay (1.04), interpreting lag (1) \\
$r$, $t_r$, $R$ & index; seconds; count & read, its time $\min(r\Delta,T)$, number of reads \\
$\bm{P}^{(r)}$ & $[0,1]^{F\times S}$ & diarizer activities visible at read $r$ (frames after $t_r-\gamma$ zeroed) \\
$\eta_r$ & $\{\wait,\eos\}$ & stop token of read $r$ \\
$e_j$, $r(j)$ & seconds; read index & source end of chunk $j$ and the read at which it is due \\
$\bm{v}_k$ & $\le8$\,s of audio & voice prompt of slot $k$, encoded into 32 BiCodec global tokens \\
\midrule
\multicolumn{3}{@{}l}{\emph{Playback alignment (\S\ref{sec:method-plem}), one speaker stream $\bm{x}$ with blocks $i=1,\dots,I$}} \\
$W_i$, $\nu_i$ & $\mathbb{N}$; seconds & target units (words or characters) of block $i$; its release time \\
$C_i$, $c$ & set, $|C_i|=8$ & candidate renderings of block $i$ (base-Talker samples at two rates) \\
$\bm{a}(c)$, $D(c)$ & semantic-token sequence; seconds & speech codes of $c$ (one per 20\,ms) and duration $0.02|\bm{a}(c)|$ \\
$\bar{\ell}(c)$, $f(c)$ & $\R_{\le0}$ (nats/token); $[-10,0]$ (nats/char) & Talker mean log-probability; negative CTC loss per character (fidelity) \\
$\mathcal{H}$, $h_i$ & 101 values in $[0,8]$\,s & backlog grid ($\delta=80$\,ms) and backlog of block $i$ \\
$g_i$ & $\mathcal{H}\times C_i\to\mathcal{H}$ & backlog transition, Eq.~\eqref{eq:backlog} \\
$L_i(h,c)$ & unit-seconds & total delay at which block $i$'s units are heard, Eq.~\eqref{eq:lat} \\
$\bm{z}$, $z_i$, $\mathcal{Z}(\bm{x})$ & $\prod_i(\mathcal{H}\times C_i)$ & playback plan, its state at block $i$, feasible plans \\
$\phi^\lambda_i$ & $\mathcal{H}\times C_i\to\R$ & node potential, Eq.~\eqref{eq:plpot}, weights $(0.5,5)$ fixed \\
$\lambda$, $\lambda^\star$ & $\R_{\ge0}$ (per unit-second) & price of delay; its value at the budget \\
$Z_\lambda(\bm{x})$ & $\R_{>0}$ & partition function \\
$\alpha_i$, $\beta_i$ & $\R_{\ge0}^{|\mathcal{H}|}$ & forward / backward messages, one entry per backlog value \\
$\mu_i$ & $[0,1]^{|\mathcal{H}|\times|C_i|}$ & posterior marginals: rows backlog values, columns candidates \\
$\Phi$, $\Phi_\mathcal{S}$ & unit-seconds & total delay of a plan; of a corpus $\mathcal{S}$ of streams \\
$b$, $W$ & seconds per unit; $\mathbb{N}$ & latency budget; target units in $\mathcal{S}$ \\
$\mathcal{D}(\lambda)$ & $\R$ & dual function, Eq.~\eqref{eq:dual} \\
$\tilde{h}_i$, $w_i$ & $\mathcal{H}$; $[0,1]^{|C_i|}$ & backlog under the base rate; posterior weights of block $i$'s candidates \\
$\pi_{\bm{\omega}}$, $\pi_0$ & distributions over token sequences & PA Talker (parameters $\bm{\omega}$) and base Talker \\
$\tau(c)$, $\rho_i$ & $\{1,\dots,35\}$; $[0,1]^{35}$ & rate token of $c$; posterior distribution over rate tokens \\
$\kappa_\text{rate}$, $\kappa_\text{kl}$ & $\R_{\ge0}$ & loss weights (1, 0.5) \\
\bottomrule
\end{tabular}
\caption{Notation. Spaces give the dimensions of every symbol and what each dimension indexes.}
\label{tab:notation}
\end{table*}

% Paper v12 appendix: only what the playback-alignment CRF and the main text rely on.

\section{Playback Alignment: Proofs}
\label{app:proofs}

\begin{figure}[t]
\centering
\begin{tikzpicture}[x=0.82cm,y=0.6cm,font=\scriptsize,
  st/.style={circle,draw,inner sep=0pt,minimum size=4.2mm},
  on/.style={st,fill=trained!25},
  blk/.style={rectangle,draw,rounded corners=1pt,inner sep=1.5pt,minimum height=3.2mm}]
% time axis with release times
\draw[->,gray] (-0.3,-3.75) -- (7.4,-3.75) node[right,black]{$t$};
\foreach \x/\l in {0/\nu_i,2.4/\nu_{i+1},4.8/\nu_{i+2}}{\draw[gray] (\x,-3.65)--(\x,-3.85) node[below,black]{$\l$};}
% backlog states (rows = backlog values)
\foreach \c/\x in {i/0,i+1/2.4,i+2/4.8}{\node[gray] at (\x,1.05) {$h_{\c}$};}
\foreach \i/\x in {0/0,1/2.4,2/4.8}{\foreach \j in {0,1,2}{\node[st] (s\i\j) at (\x,-\j) {};}}
\node[left,gray] at (-0.25,0) {0\,s};\node[left,gray] at (-0.25,-1) {0.4\,s};\node[left,gray] at (-0.25,-2) {0.8\,s};
% candidate edges from h_i = 0
\draw[->,thick,trained] (s00) -- node[above,sloped]{$c$: short} (s10);
\draw[->,thick,frozen,dashed] (s00) -- node[below,sloped,pos=0.6]{long} (s12);
\draw[->,thin,gray] (s10) -- (s20); \draw[->,thin,gray] (s10) -- (s21);
\draw[->,thin,gray] (s12) -- (s21); \draw[->,thin,gray] (s12) -- (s22);
% backward message
\draw[->,densely dotted,thick,spkc] (5.6,-2.4) to[bend left=12] node[below,pos=0.45]{backward message $\beta_{i+1}(h)$} (0.6,-2.55);
\node[right,align=left] at (5.15,-0.35) {$\phi_i(h,c)=$\\$0.5\bar\ell(c)+5f(c)$\\$-\lambda L_i(h,c)$};
\end{tikzpicture}
\caption{Playback alignment for one speaker stream. States are the backlog $h_i$ (time block $i$ waits for earlier
speech); an edge is a candidate rendering $c$ of block $i$, whose duration sets the next backlog
(Eq.~\eqref{eq:backlog}). A long candidate delays every later block; backward messages bring that cost into the choice
for block $i$.}
\label{fig:pa}
\end{figure}


Figure~\ref{fig:pa} shows the state space of Eq.~\eqref{eq:crf} for one stream.

Throughout, fix $\lambda$ and write $\phi_i=\phi^\lambda_i$. A feasible plan is determined by its candidates: given
$c_{1:I}$, Eq.~\eqref{eq:backlog} fixes every backlog from $h_1=0$. We write $h_j(c_{<j})$ for the backlog that
$c_{<j}$ leads to, and $e(\bm{z})=\exp\sum_i\phi_i(z_i)$ for the unnormalised weight of a plan.

\begin{proposition}[Exact inference]
\label{prop:fb}
For the CRF of Eq.~\eqref{eq:crf}, the messages of Eq.~\eqref{eq:msg} satisfy the recursions of Eq.~\eqref{eq:fb};
for every $i$, $Z_\lambda=\sum_h\alpha_i(h)\beta_i(h)$, in particular $Z_\lambda=\beta_1(0)$; and the marginals are given
by Eq.~\eqref{eq:marg}. The cost is $O(I\,|\mathcal{H}|\max_i|C_i|)$.
\end{proposition}

\begin{proof}
\emph{Factorisation.} By Eq.~\eqref{eq:backlog}, the backlogs after block $i$ are computed from $h_i$ and
$c_i,c_{i+1},\dots$ alone: $c_{<i}$ influences them only through $h_i$. Hence the weight of a plan splits at any block,
\[
e(\bm{z})=\Big(\textstyle\prod_{j<i}e^{\phi_j(h_j,c_j)}\Big)\Big(\textstyle\prod_{j\ge i}e^{\phi_j(h_j,c_j)}\Big),
\]
a prefix factor that ends at $h_i$ times a suffix factor that starts at $h_i$ and depends on the prefix only through
$h_i$.
\emph{Partition function.} Summing $e(\bm{z})$ over all feasible plans and grouping them by the value $h$ of $h_i$, the
prefix sums give $\alpha_i(h)$ and the suffix sums $\beta_i(h)$ (Eq.~\eqref{eq:msg}), so $Z_\lambda=\sum_h\alpha_i(h)\beta_i(h)$.
At $i=1$, $\alpha_1(h)=\mathds{1}[h=0]$ gives $Z_\lambda=\beta_1(0)$.
\emph{Recursions.} A prefix ending at $h_{i+1}=h'$ is a prefix ending at some $h_i=h$ followed by a candidate $c$ with
$g_i(h,c)=h'$; its weight is the old weight times $e^{\phi_i(h,c)}$, which gives $\alpha_{i+1}$. A suffix starting at
$h_i=h$ is a candidate $c$ followed by a suffix starting at $g_i(h,c)$, which gives $\beta_i$.
\emph{Marginals.} The plans with $z_i=(h,c)$ are a prefix ending at $h$, the state $(h,c)$ and a suffix starting at
$g_i(h,c)$, with total weight $\alpha_i(h)e^{\phi_i(h,c)}\beta_{i+1}(g_i(h,c))$; dividing by $Z_\lambda$ gives
Eq.~\eqref{eq:marg}, and $\mathbb{E}_{p_\lambda}[\Phi]=\sum_i\mathbb{E}[L_i(z_i)]$ by linearity.
\emph{Cost.} Each recursion step visits every pair $(h,c)\in\mathcal{H}\times C_i$ once.
\end{proof}

\begin{proposition}[Latency budget]
\label{prop:pr}
Assume some distribution meets the budget of Eq.~\eqref{eq:pr} strictly (it suffices that the minimum-delay plans do).
Then the solution of Eq.~\eqref{eq:pr} is $q^\star=\prod_{\bm{x}\in\mathcal{S}}p_{\lambda^\star}(\cdot\mid\bm{x})$, where
$\lambda^\star$ maximises the dual of Eq.~\eqref{eq:dual} over $\lambda\ge0$; the dual has the derivatives of
Eq.~\eqref{eq:dualderiv}, and $\lambda^\star=0$ if $\mathbb{E}_{p_0}[\Phi_\mathcal{S}]\le bW$, otherwise
$\mathbb{E}_{p_{\lambda^\star}}[\Phi_\mathcal{S}]=bW$.
\end{proposition}

\begin{proof}
\emph{Inner minimisation.} For $\lambda\ge0$, define the Gibbs distribution
$q_\lambda(\bm{z})=p_0(\bm{z})e^{-\lambda\Phi_\mathcal{S}(\bm{z})}/M_\lambda$ with $M_\lambda=\mathbb{E}_{p_0}[e^{-\lambda\Phi_\mathcal{S}}]$.
Rearranging the Lagrangian,
\[
\begin{aligned}
&\mathrm{KL}(q\Vert p_0)+\lambda\,\mathbb{E}_q[\Phi_\mathcal{S}]\\
&\quad=\textstyle\sum_{\bm{z}}q\log\big(q/(p_0e^{-\lambda\Phi_\mathcal{S}})\big)\\
&\quad=\mathrm{KL}(q\Vert q_\lambda)-\log M_\lambda ,
\end{aligned}
\]
which is minimised, uniquely, by $q=q_\lambda$.
\emph{The minimiser is the CRF.} $p_0$ is a product over streams of Eq.~\eqref{eq:crf} at $\lambda=0$, and
$\Phi_\mathcal{S}$ is a sum of per-block delays, so
$p_0(\bm{z})e^{-\lambda\Phi_\mathcal{S}(\bm{z})}=\prod_{\bm{x}}\mathds{1}[\bm{z}_{\bm{x}}\in\mathcal{Z}(\bm{x})]\exp\big(\sum_i\phi^0_i-\lambda L_i\big)/Z_0(\bm{x})$;
since $\phi^0_i-\lambda L_i=\phi^\lambda_i$, normalising gives $q_\lambda=\prod_{\bm{x}}p_\lambda(\cdot\mid\bm{x})$ and
$M_\lambda=\prod_{\bm{x}}Z_\lambda(\bm{x})/Z_0(\bm{x})$.
\emph{Dual and its derivatives.} Substituting $q_\lambda$ into the Lagrangian gives
$\mathcal{D}(\lambda)=-\log M_\lambda-\lambda bW$, i.e.\ Eq.~\eqref{eq:dual}. For one stream,
$\partial_\lambda\log Z_\lambda=\sum_{\bm{z}}(-\Phi(\bm{z}))e^{\sum_i\phi^\lambda_i}/Z_\lambda=-\mathbb{E}_{p_\lambda}[\Phi]$, and
differentiating once more, $\partial^2_\lambda\log Z_\lambda=\mathbb{E}_{p_\lambda}[\Phi^2]-\mathbb{E}_{p_\lambda}[\Phi]^2=\mathrm{Var}_{p_\lambda}[\Phi]$.
Streams are independent under $q_\lambda$, so expectations and variances add, which gives Eq.~\eqref{eq:dualderiv}.
\emph{Optimality.} Eq.~\eqref{eq:pr} minimises a convex function of $q$ under one linear constraint, and a strictly
feasible point exists, so strong duality holds and $q^\star=q_{\lambda^\star}$ with $\lambda^\star=\arg\max_{\lambda\ge0}\mathcal{D}$.
As $\mathcal{D}''\le0$, $\mathcal{D}'(\lambda)=\mathbb{E}_{p_\lambda}[\Phi_\mathcal{S}]-bW$ is non-increasing: if it is
$\le0$ at $\lambda=0$ the maximum is at $\lambda^\star=0$, otherwise at its root, where the budget holds with equality
(complementary slackness). A non-increasing function has its root found by bisection.
\end{proof}

\begin{proposition}[Posterior distillation]
\label{prop:distill}
Let $\pi$ be a distribution over speech-code sequences, $C$ a set of candidates with distinct code sequences
$\bm{a}(c)$, $\pi(C)=\sum_{c\in C}\pi(\bm{a}(c))>0$ and $\tilde{\pi}(c)=\pi(\bm{a}(c))/\pi(C)$. For any distribution $w$
on $C$,
$-\sum_{c\in C}w(c)\log\pi(\bm{a}(c))=\mathrm{KL}(w\Vert\tilde{\pi})+\mathrm{H}(w)-\log\pi(C)$.
\end{proposition}

\begin{proof}
Write $\log\pi(\bm{a}(c))=\log\tilde{\pi}(c)+\log\pi(C)$ and use $\sum_cw(c)=1$:
$-\sum_cw(c)\log\pi(\bm{a}(c))=-\sum_cw(c)\log\tilde{\pi}(c)-\log\pi(C)$, and
$-\sum_cw(c)\log\tilde{\pi}(c)=\mathrm{KL}(w\Vert\tilde{\pi})+\mathrm{H}(w)$.
Since $\mathrm{H}(w)$ does not depend on $\pi$, minimising the weighted cross-entropy over $\pi$ moves $\tilde{\pi}$
towards the CRF posterior $w$ and raises the mass $\pi(C)$ on renderings of the kind the CRF scored.
\end{proof}

\section{Playback Alignment: Candidates, Training and Base Rate}
\label{app:plem}

\begin{figure}[h]
\small
\resizebox{\columnwidth}{!}{\begin{tabular}{@{}l@{}}
\toprule
\textbf{Algorithm 1}: training the PA Talker (one direction) \\
\midrule
1: decode training meetings with a streaming Thinker; merge each read's blocks per speaker \\
2: play every stream at the base rate; record $\nu_i$ and the backlog $\tilde{h}_i$ \\
3: for every block, sample 8 candidates from the base Talker; compute $D$, $\bar{\ell}$, $f$ \\
4: bisect $\lambda$ until $\mathbb{E}_{p_\lambda}[\Phi_\mathcal{S}]=bW$ (Eq.~\eqref{eq:marg}; Proposition~\ref{prop:pr}) \\
5: run the backward recursion at $\lambda^\star$; weight candidates by $w_i(c)$ (Eq.~\eqref{eq:w}) \\
6: fine-tune the Talker on $\mathcal{J}$ (Eq.~\eqref{eq:distill}), candidates with $w_i(c)\ge0.01$ \\
\textbf{Test}: backlog token of the actual queue; 2 samples per block; keep the higher $f$ \\
\bottomrule
\end{tabular}}
\caption{Training and decoding of the playback layer. No CRF runs at test time.}
\label{alg:pa}
\end{figure}

\paragraph{Candidates and fidelity.} For every block the base Talker samples 8 candidates: 4 at the rate the
base rate gives and 4 at two rate steps faster. $f(c)=-\min(\text{CTC loss per character},10)$ under an
independent wav2vec2 CTC recogniser~\citep{baevski2020wav2vec} of the target language. The forward and backward
recursions agree with brute-force enumeration on small streams.

\paragraph{Objective.} The Talker, conditioned on a backlog token (0.5\,s bins, capped at 3.5\,s), minimises
Eq.~\eqref{eq:distill}: the weighted cross-entropy of the candidates' speech codes, plus \num{v10:c:plem:wrate} times the
cross-entropy of its rate token (\num{v10:c:plem:rates} rates) against the posterior's rate marginal, plus
\num{v10:c:plem:wkl} times the KL divergence from the frozen base Talker on the semantic codes and the end token. All
Talker parameters are trained.

\paragraph{Data.} Zh$\rightarrow$En: \num{v10:c:plem:meet:z} AMI training meetings, decoded by a separately trained
single-direction streaming Thinker, \num{v10:c:plem:blocks:z} blocks, \num{v10:c:plem:rows:z} weighted candidate rows,
$\lambda^\star=$ \num{v10:c:plem:lam:z}, \num{v10:c:plem:steps:z} steps. En$\rightarrow$Zh: the same recipe,
\num{v10:c:plem:rows:e} rows, $\lambda^\star=$ \num{v10:c:plem:lam:e}; of its checkpoints at iterations 400, 800 and
1{,}200 a rule on five development meetings (highest ASR-BLEU among checkpoints within 0.1\,s of the last one's speech
latency) chose iteration \num{v10:c:plem:it:e}. The budget is the base-rate path's own delay minus
\num{v10:c:plem:budgetoff}\,s.

\paragraph{Base rate and ablation rows.} The base rate of a block is its source pace plus three rate steps per second
of backlog (at most six), shifted by two steps. In Table~\ref{tab:main}, R1 merges the blocks of one read, R2 speaks
them with the base Talker at the base rate, R3 replaces the base Talker by the PA Talker with one sample, and R4
(\sys{}) keeps the better of two samples by $f$; R3 $-$ R2 isolates the CRF-distilled Talker and R4 $-$ R3 the
selection. Differences between rows are descriptive.

\section{Streaming Thinker: Model and Training}
\label{app:model}
\label{app:train}
\label{app:causal}
\label{app:prompt}

\paragraph{Speaker tokens.} The diarizer's posteriors $\bm{P}\in[0,1]^{F\times S}$ (80\,ms frames, $S=8$ slots) are
resampled to the acoustic token centres, $\tilde{\bm{P}}=\mathrm{Interp}(\bm{P};\bar{t}_{1:N})$, zero-padded to $MK$
rows ($M=\lceil N/K\rceil$), and every block of $K=5$ rows becomes one speaker token,
\begin{equation}
\begin{aligned}
\bm{e}_j&=\textstyle\sum_{r=1}^{K}\bm{W}^{(r)}\tilde{\bm{p}}_{(j-1)K+r}+\bm{b},\\
\bm{s}_j&=\bm{W}_2\,\mathrm{GELU}\big(\bm{W}_1\mathrm{GELU}(\bm{e}_j)+\bm{b}_1\big)+\bm{b}_2,
\end{aligned}
\label{eq:conv}
\end{equation}
with a hidden size of 1{,}024; $\bm{W}_2,\bm{b}_2$ start at zero, so training starts from the pretrained Thinker.
Per meeting, reference speakers $\rho$ are mapped to slots by the injective map $\sigma^\star$ that maximises the
number of frames in which $\rho$ speaks and slot $\sigma(\rho)$ is active ($P>0.5$); an utterance by $\rho$ is
labelled $\sigma^\star(\rho)$. Hesitations have no target segment. The instruction asks for
\texttt{<|spk k|><|t a|>}translation\texttt{<|t b|>} per utterance, $k$ the slot, $a,b$ in 80\,ms steps, utterances
in order of start.

\paragraph{Deadlines.} An utterance's reference translation is split at punctuation into chunks $c_1,\dots,c_J$; chunk
$j$ covers the fraction $\varphi_j=\sum_{j'\le j}|c_{j'}|/\sum_{j'}|c_{j'}|$ of the utterance and ends at the end of
source word $\lceil\varphi_j|\bm{w}|\rceil$ of the time-aligned words $\bm{w}$, or at $a+\varphi_j(b-a)$ where no
word timings exist (the re-voiced Chinese source).

\paragraph{Voices and playback layout.} Slot $k$'s solo frames are
\begin{equation}
m_{f,k}=\mathds{1}\big[P_{f,k}>0.5\big]\textstyle\prod_{k'\neq k}\mathds{1}\big[P_{f,k'}<0.5\big].
\label{eq:alone}
\end{equation}
Within a speaker's stream, segment $n$ of duration $D_n$ starts at
\begin{equation}
\pi_n=\max\big(\zeta_n,\;\pi_{n-1}+D_{n-1}\big),
\label{eq:layout}
\end{equation}
with $\zeta_n$ its release time; this is the queue of Eq.~\eqref{eq:backlog} before discretisation.

\paragraph{Causality.} Streaming runs read the posteriors of one low-latency diarizer run over the whole session, with
every frame later than $t-\gamma$ zeroed at a read at time $t$. For three windows of every test meeting we truncated the
audio at a read time and re-ran the diarizer from scratch: over \num{v6:caus:cuts} cuts the largest absolute posterior
difference was \num{v6:caus:maxabs} (fp16 precision) and every differing $>0.5$ decision was an fp16 rounding tie
within $5\cdot10^{-4}$ of 0.5. The audio encoder encodes the heard prefix at every read.

\paragraph{Training.} The deadline rows of the two directions are interleaved into \num{v10:c:j:rows} rows such that
every global batch of 64 holds 32 Zh$\rightarrow$En and 32 En$\rightarrow$Zh rows; learning rate cosine from $10^{-4}$
to $10^{-5}$ over \num{v10:c:j:s1steps} steps (100 warm-up), LoRA, projector and new rows as in Table~\ref{tab:hyper},
8 H200 GPUs, micro-batch 8. There is no checkpoint selection: the Thinker is the last step. Reads with nothing due are
kept with probability 0.5; whole-window (20\%) and single-speaker (15\%) rows are mixed in.

\begin{table}[h]
\centering\small
\setlength{\tabcolsep}{3pt}
\resizebox{\columnwidth}{!}{\begin{tabular}{lcc}
\toprule
 & \sys{} (offline) & \sys{}-Simul \\
\midrule
Initialisation & Qwen3-Omni Thinker & \sys{} (offline), merged \\
Directions & both & one model each \\
Diarizer preset & offline & low latency (1.04\,s) \\
Rows & 67{,}795 & 128{,}879 / 160{,}957 \\
Steps & 2{,}000 & 3{,}000 \\
Batch (global / micro) & 64 / 8 & 64 / 8 \\
Learning rate & \multicolumn{2}{c}{$10^{-4}\to10^{-5}$, cosine, 50 warm-up} \\
LoRA & \multicolumn{2}{c}{$r=32$, $\alpha=64$, q/k/v/o, 48 layers} \\
Loss weights $\omega$ & \multicolumn{2}{c}{speaker 2, time 1.5, other 1} \\
Posterior perturbation & \multicolumn{2}{c}{$p=0.2$, $\mathcal{N}(0,0.1^2)$, shift $\le1$ frame} \\
Read interval / lag & -- & $\Delta=1$\,s / $\epsilon=1$\,s \\
\bottomrule
\end{tabular}}
\caption{Training hyperparameters (\sys{}-Simul: En$\rightarrow$Zh / Zh$\rightarrow$En). The streaming Thinker of
\sys{} is one bidirectional model trained like \sys{}-Simul on the rows of both directions.}
\label{tab:hyper}
\end{table}

\section{Data}
\label{app:data}

\input{generated/tab_data}

\paragraph{Translations.} Every AMI utterance (IHM segmentation, official split) is translated by three MT systems
(Hy-MT2-30B-A3B, NiuTrans LMT-60-4B, MiLMMT-46-4B); one candidate is selected by reference-free quality estimation
(XCOMET-XL~\citep{guerreiro2024xcomet} and an XLM-R-XL estimator) and chrF agreement, synthesised in the speaker's
voice and checked for audibility and intelligibility. A window enters the data only if all its non-hesitation
utterances have a translation whose speech decodes. Reference windows group utterances that overlap or touch, up to
20\,s.

\paragraph{Chinese AMI meetings.}\label{app:zhdata} Every utterance's Chinese translation is re-spoken by the UniSS Talker with one voice
prompt per speaker, trimmed, set to the speaker's median loudness in the English recording, placed at the utterance's
original start and time-stretched towards the English duration within $[0.8,1.5]$; a speaker's next clip waits until
the previous one ends. Clips failing an intelligibility check (CER $>0.3$) are left out of the meeting, and 13 training
meetings whose diarization error stayed high were dropped (none in test).

\paragraph{AliMeeting references.} The main results use DiariST's human English translations~\citep{yang2024diarist};
the overlap study (Table~\ref{tab:ovmain8}) uses machine translations (Hy-MT2-30B) of the human transcripts, without a human check.

\section{Metrics and Latency}
\label{app:metrics}
\label{app:chunk}

\paragraph{tcpBLEU.} Reference units are a speaker's scored utterances chained while less than 1\,s apart; output
slots are matched to speakers per meeting by chrF over the meeting's text; each output segment, on the meeting clock and
widened by a 2\,s collar, goes to the annotated utterance it overlaps most: to that utterance's unit if it belongs to the
segment's matched speaker, otherwise to the matched speaker's own unit it overlaps most, and otherwise it is an
insertion. tcpBLEU is corpus BLEU over (unit hypothesis, unit reference) pairs plus (insertion, empty) pairs.

\paragraph{Chunk protocol.} The low-latency diarizer emits frames in blocks of nine; frame $f$ (80\,ms) is known at
$t^{\mathrm{rel}}_f=(9\lfloor f/9\rfloor+13)\cdot 80\,\text{ms}\le t_f+\gamma$, $\gamma=1.04$\,s. A window's opening is
released at $t^{\mathrm{rel}}$ of its first active frame and its end at $t^{\mathrm{rel}}$ of the last frame of the
closing gap; commits before the opening are delayed to it, and those of the last read to the release of the end.

\paragraph{Speech LAAL.} For a reference unit with source duration $|X|$ and $|Y|$ heard target words (words or Chinese
characters) with delays $d_1\le\dots\le d_{|Y|}$ from the unit's start,
$\mathrm{LAAL}=\frac{1}{k^\ast}\sum_{k=1}^{k^\ast}\big(d_k-(k-1)\frac{|X|}{\max(|Y|,|Y^\ast|)}\big)$, with $|Y^\ast|$ the
reference length and $k^\ast$ the first $k$ with $d_k\ge|X|$~\citep{papi2022over}. Block $i$ of a stream starts playing
at $\nu_i+h_i$. LAAL$_\text{block}$ gives every word of the block $d_k=\nu_i+h_i$ (relative to the unit's start);
LAAL$_\text{spread}$ spreads its $W_i$ words evenly over the block's duration, as Eq.~\eqref{eq:lat}. Speech start is the
first heard word of a unit; EndOffset runs from the unit's end to the end of its last heard word.

\paragraph{Where the delay comes from.} Table~\ref{tab:decomp} splits EndOffset into the model's lag, the detector's
wait, the wait behind earlier speech of the same speaker (the backlog of Eq.~\eqref{eq:backlog}) and the duration of the
last block. The base Talker produces more speech than the source holds (load above one), so each stream accumulates a
queue; the playback layer removes most of that queue and shortens the last block, the two parts of its gain in
EndOffset.

\begin{table}[t]
\centering\small
\setlength{\tabcolsep}{2.5pt}
\resizebox{\columnwidth}{!}{\begin{tabular}{llrrrrrrr}
\toprule
Set & Sys. & start & model & detector & queue & last block & EndOff. & load \\
\midrule
dev Zh$\rightarrow$En & B & 2.95 & 0.74 & 0.61 & 0.64 & 2.16 & 4.14 & 1.41 \\
 & \sys{} & 2.90 & 0.75 & 0.61 & 0.12 & 1.62 & 3.10 & 0.94 \\
\midrule
dev En$\rightarrow$Zh & B & 2.36 & 0.76 & 0.55 & 0.59 & 1.54 & 3.45 & 1.33 \\
 & \sys{} & 2.32 & 0.78 & 0.55 & 0.14 & 1.37 & 2.83 & 1.00 \\
\midrule
AMI Zh$\rightarrow$En & B & 2.99 & 0.67 & 0.66 & 0.65 & 2.10 & 4.08 & 1.37 \\
 & \sys{} & 2.96 & 0.68 & 0.66 & 0.14 & 1.64 & 3.13 & 0.94 \\
\midrule
AliMeeting & B & 3.05 & 0.88 & 0.40 & 1.33 & 1.96 & 4.57 & 1.56 \\
 & \sys{} & 2.92 & 0.95 & 0.40 & 0.33 & 1.55 & 3.23 & 1.05 \\
\midrule
AMI En$\rightarrow$Zh & B & 2.43 & 0.69 & 0.64 & 0.54 & 1.61 & 3.48 & 1.28 \\
 & \sys{} & 2.39 & 0.71 & 0.64 & 0.14 & 1.42 & 2.91 & 0.99 \\
\bottomrule
\end{tabular}}
\caption{Where the speech delay of a source unit comes from (chunk protocol, mean over units and two TTS seeds, seconds). EndOffset = model lag (last commit after the unit's end) + detector wait (online window release) + queue wait (earlier speech of the same stream) + duration of the last block. start: speech start offset from the unit's start. load: synthesised speech per second of source speech. \sys{}: the deadline Thinker's text with read-level merging and the PA Talker.}
\label{tab:decomp}
\end{table}


\section{Further Results of the Playback Layer}
\label{app:sig10}
\label{app:postreview}

\paragraph{Significance.} Table~\ref{tab:sig} lists the paired BLEU tests of \sys{} against B. The speech of the two
systems is compared on the same text; Table~\ref{tab:speech} gives both speech-LAAL definitions, speech start,
EndOffset, speaking rate, UTMOS~\citep{saeki2022utmos}, speaker similarity (WavLM, against the Talker's voice prompt)
and pauses, with paired latency differences and their intervals. Every latency definition falls; the speaking rate rises
by \num{v11:pr:test:ami:e2-b:ratepct} / \num{v11:pr:test:ali:e2-b:ratepct} / \num{v11:pr:test:en:e2-b:ratepct}\%
(AMI Zh$\rightarrow$En / AliMeeting / En$\rightarrow$Zh); UTMOS, paired per stream, rises by
\num{v10:pr:test:ami:e2-b:utmos:diff} / \num{v10:pr:test:ali:e2-b:utmos:diff} / \num{v10:pr:test:en:e2-b:utmos:diff};
speaker similarity is unchanged; pauses become two to three times as frequent.

\begin{table}[t]
\centering\small
\setlength{\tabcolsep}{2.5pt}
\resizebox{\columnwidth}{!}{\begin{tabular}{llrrrr}
\toprule
Set & Comparison & $\Delta$ & 95\% interval & $p$ & $p_\text{Holm}$ \\
\midrule
 & ASR-BLEU: \sys{} $-$ B & +2.28 & [1.57, 3.04] & $<$0.001 & $<$0.001 \\
\midrule
 & ASR-BLEU: \sys{} $-$ B & +1.35 & [1.09, 2.28] & $<$0.001 & $<$0.001 \\
\midrule
 & ASR-BLEU: \sys{} $-$ B & +0.58 & [0.43, 0.74] & $<$0.001 & $<$0.001 \\
\bottomrule
\end{tabular}}
\caption{Paired tests over meetings on the test sets (16 AMI, 20 AliMeeting meetings): exact two-sided sign-flip permutation $p$ over all $2^M$ flips, 95\% meeting-bootstrap interval (10{,}000 resamples), Holm over the twelve BLEU tests of the analysis plan. Speech pools the two TTS seeds per meeting. B: deadline baseline; \sys{}: its text with rl + PA Talker. The text of \sys{} equals B's by construction. The intervals do not include decoder run-to-run noise; training-seed variance is not covered.}
\label{tab:sig}
\end{table}

\begin{table*}[t]
\centering\small
\setlength{\tabcolsep}{3.5pt}
\begin{tabular}{llrrrrrrrrr}
\toprule
 & & \multicolumn{4}{c}{speech latency (s)} & & \multicolumn{2}{c}{voice} & \multicolumn{2}{c}{pauses / min} \\
\cmidrule(lr){3-6}\cmidrule(lr){8-9}\cmidrule(lr){10-11}
Set & Sys. & LAAL$_\text{block}$ & LAAL$_\text{spread}$ & start & EndOffset & rate & UTMOS & SIM & $\ge$0.5\,s & mid-sent. \\
\midrule
AMI Zh$\rightarrow$En & B & 3.12 & 4.01 & 2.99 & 4.08 & 2.40 & 1.80 & 0.827 & 3.3 & 0.8 \\
 & \sys{} & 3.03 & 3.61 & 2.96 & 3.13 & 3.49 & 1.94 & 0.827 & 8.3 & 2.8 \\
 & $\bias=0$ & 3.39 & 3.99 & 3.39 & 3.42 & 3.49 & -- & -- & 6.0 & 0.8 \\
\midrule
AliMeeting & B & 3.28 & 4.28 & 3.05 & 4.57 & 2.40 & 3.48 & 0.873 & 4.4 & 0.9 \\
 & \sys{} & 2.95 & 3.65 & 2.92 & 3.23 & 3.56 & 3.72 & 0.873 & 11.8 & 3.9 \\
 & $\bias=0$ & 4.01 & 4.85 & 4.15 & 4.23 & 3.54 & -- & -- & 8.3 & 1.2 \\
\midrule
AMI En$\rightarrow$Zh & B & 2.54 & 3.21 & 2.43 & 3.48 & 3.90 & 3.25 & 0.833 & 4.6 & -- \\
 & \sys{} & 2.43 & 2.95 & 2.39 & 2.91 & 5.04 & 3.40 & 0.831 & 9.4 & -- \\
 & $\bias=0$ & 2.79 & 3.34 & 2.81 & 3.21 & 5.02 & -- & -- & 7.8 & -- \\
\bottomrule
\end{tabular}
\caption{Speech of the streaming systems on the test sets. Latency: chunk protocol, mean of two TTS seeds; block-start LAAL charges every word of a block when the block starts playing, LAAL$_\text{spread}$ spreads them over its playback; start: first heard word of a unit. Rate: English words (Zh$\rightarrow$En) or characters (En$\rightarrow$Zh) per second of synthesised speech. UTMOS and SIM (WavLM speaker similarity to the Talker's voice prompt, not to the speaker's own recording): one TTS seed, all streams of at least 0.5\,s. Pauses: silences of at least 0.5\,s per minute of speech; mid-sentence: after a block that does not end a sentence (English output only).}
\label{tab:speech}
\end{table*}

\begin{table*}[t]
\centering\small
\setlength{\tabcolsep}{3pt}
\resizebox{\textwidth}{!}{\begin{tabular}{llllll}
\toprule
Set & Comparison & LAAL$_\text{block}$ & LAAL$_\text{spread}$ & start & EndOffset \\
\midrule
 & \sys{} $-$ B & $-$0.09 [$-$0.10, $-$0.08] & $-$0.40 [$-$0.42, $-$0.38] & $-$0.03 [$-$0.04, $-$0.02] & $-$0.96 [$-$1.01, $-$0.91] \\
\midrule
 & \sys{} $-$ B & $-$0.33 [$-$0.40, $-$0.28] & $-$0.63 [$-$0.69, $-$0.57] & $-$0.13 [$-$0.17, $-$0.10] & $-$1.34 [$-$1.45, $-$1.23] \\
\midrule
 & \sys{} $-$ B & $-$0.11 [$-$0.13, $-$0.09] & $-$0.26 [$-$0.28, $-$0.24] & $-$0.04 [$-$0.06, $-$0.03] & $-$0.58 [$-$0.62, $-$0.54] \\
\bottomrule
\end{tabular}}
\caption{Paired speech-latency differences on the test sets (seconds): difference and 95\% meeting-bootstrap interval (10{,}000 resamples, two TTS seeds pooled per meeting). Exact sign-flip $p$ values are given in the text where used; no multiple-test correction.}
\label{tab:speech2}
\end{table*}


\ifAliEval
\subsection{Locked test on unused AliMeeting meetings}
\label{app:alieval}
The eight AliMeeting evaluation meetings (\num{v11:ae:hours}\,h) were never used for training, development or any
earlier test. Windows, diarizer runs and scoring follow the AliMeeting test set, except that the references are
DiariST's human English translations and no window is filtered by its reference. The analysis plan fixed the systems
and three confirmatory hypotheses (\sys{} $-$ B on ASR-BLEU, LAAL$_\text{block}$ and EndOffset, Holm over three)
before decoding; each system was decoded once (Table~\ref{tab:alieval}).
\begin{table*}[t]
\centering\small
\setlength{\tabcolsep}{3.5pt}
\begin{tabular}{lrrrrrrrrrr}
\toprule
 & text & ASR- & text & \multicolumn{4}{c}{speech latency (s)} & rate & pauses & \\
\cmidrule(lr){5-8}
System & tcpBLEU & BLEU & LAAL & LAAL$_\text{block}$ & LAAL$_\text{spread}$ & start & EndOff. & (w/s) & /min & UTMOS \\
\midrule
B & 14.58 & 10.66 & 2.87 & 3.39 & 4.50 & 3.23 & 4.79 & 2.36 & 4.3 & 3.50 \\
\sys{} & 14.58 & 12.89 & 2.87 & 3.02 & 3.81 & 3.09 & 3.35 & 3.50 & 11.1 & 3.77 \\
\bottomrule
\end{tabular}
\caption{Locked test on the eight unused AliMeeting evaluation meetings (4.2\,h, Zh$\rightarrow$En, DiariST human references), each system decoded once; ASR-BLEU and speech metrics averaged over two TTS seeds (UTMOS: seed 1), latencies under the chunk protocol without computation. Confirmatory tests of \sys{} $-$ B (Holm over three; exact sign-flip over 8 meetings, minimum $p$ = 2/256; 95\% meeting-bootstrap interval): ASR-BLEU +2.23 [+0.75, +3.37], $p_\text{Holm}$ 0.023; LAAL$_\text{block}$ $-$0.37 [$-$0.41, $-$0.33], $p_\text{Holm}$ 0.023; EndOffset $-$1.43 [$-$1.51, $-$1.32], $p_\text{Holm}$ 0.023.}
\label{tab:alieval}
\end{table*}

\fi

\ifMetricsV
\subsection{A second recogniser and COMET}
\label{app:metricsv11}
Candidate selection uses a CTC recogniser (wav2vec~2.0), and ASR-BLEU uses Whisper-large-v3 or Paraformer, so part of
the gain could be speech that suits recognisers rather than listeners. We rescored the existing test outputs with
definitions fixed before computing: the speech of both TTS seeds transcribed by a second recogniser of a different
family (English: Parakeet-TDT-0.6B-v2; Chinese: Whisper-large-v3, whose repetition loops and video-outro hallucinations
are removed by one rule for all systems), and COMET~\citep{rei2022comet} per reference unit (Table~\ref{tab:metricsv11}).
Against B, \sys{} gains on the second recogniser by \num{v11:asr2:gain:ami} / \num{v11:asr2:gain:en} /
\num{v11:asr2:gain:ali} and in speech COMET on all three sets. The En$\rightarrow$Zh second-recogniser gain is larger
than with Paraformer because Whisper fails more often on B's audio (loops or hallucinations in 16\% of B's utterances
against 6\% of \sys{}'s), so only its sign is informative.
\begin{table*}[t]
\centering\small
\setlength{\tabcolsep}{4pt}
\begin{tabular}{lrrrrrrrrr}
\toprule
 & \multicolumn{2}{c}{ASR-BLEU} & \multicolumn{3}{c}{ASR-BLEU, second recogniser} & \multicolumn{3}{c}{COMET of speech} & COMET \\
\cmidrule(lr){2-3}\cmidrule(lr){4-6}\cmidrule(lr){7-9}
Test set & B & \sys{} & B & \sys{} & $\Delta$ & B & \sys{} & $\Delta$ & of text \\
\midrule
AMI Zh$\rightarrow$En & 18.07 & 20.35 & 19.30 & 20.61 & +1.32$^{*}$ & 73.86 & 76.44 & +2.58$^{*}$ & 79.42 \\
AMI En$\rightarrow$Zh & 31.65 & 32.23 & 28.13 & 31.87 & +3.74$^{*}$ & 73.16 & 73.50 & +0.34$^{*}$ & 78.81 \\
AliMeeting & 16.73 & 18.08 & 17.60 & 18.54 & +0.93$^{*}$ & 64.97 & 66.44 & +1.47$^{*}$ & 71.49 \\
\bottomrule
\end{tabular}
\caption{The speech-layer gain under a second recogniser (English: Parakeet-TDT-0.6B-v2, Chinese: Whisper-large-v3; neither is the CTC recogniser used for selection; repetition loops and video-outro hallucinations of the second recogniser removed by one fixed rule for all systems) and under COMET (wmt22-comet-da, $\times100$, per reference unit, empty hypotheses scored). ASR-BLEU: the paper's recognisers (Whisper-large-v3 / Paraformer). Both TTS seeds pooled; B and \sys{} share their text. $\Delta$: \sys{} $-$ B, meeting sign-flip test, Holm over the test family: $^{*}$ $p_\text{Holm}<0.05$, $^{\circ}$ not significant.}
\label{tab:metricsv11}
\end{table*}

\fi

\subsection{External system: SeamlessStreaming}
\label{app:e7}
SeamlessStreaming~\citep{seamless2023} (official agent and default settings, inference only) reads each of our online
test windows; each output segment is attributed to a diarizer slot by causal relabelling, and text, speech and latency
are scored with our tools under the chunk protocol. Nothing was trained or tuned, and each test set was decoded once.
It is a zero-shot reference for a public single-speaker SimulS2ST system on overlapping meetings, not a like-for-like
baseline (Table~\ref{tab:e7}): its quality is far below ours, most on overlapped units, and LAAL compares systems of
similar quality only.
\begin{table}[t]
\centering\small
\setlength{\tabcolsep}{2.5pt}
\resizebox{\columnwidth}{!}{\begin{tabular}{llrrrrr}
\toprule
Set & System & tcpBLEU & ov & ASR-BLEU & LAAL$_\text{spread}$ & EndOff. \\
\midrule
AMI Zh$\rightarrow$En & \sys{} & 25.02 & 19.97 & 20.35 & 3.61 & 3.13 \\
 & SeamlessStreaming & 8.69 & 3.07 & 8.48 & 3.29 & 2.53 \\
\midrule
AliMeeting & \sys{} & 21.43 & 16.91 & 18.08 & 3.65 & 3.23 \\
 & SeamlessStreaming & 10.41 & 4.42 & 9.42 & 3.96 & 2.72 \\
\midrule
AMI En$\rightarrow$Zh & \sys{} & 41.15 & 36.16 & 32.23 & 2.95 & 2.91 \\
 & SeamlessStreaming & 13.28 & 5.86 & 11.31 & 3.19 & 3.04 \\
\bottomrule
\end{tabular}}
\caption{External zero-shot single-speaker S2ST reference on our online test windows (no speaker attribution of its own: causal diarizer relabelling; one deterministic run; chunk protocol, seconds). LAAL compares systems of similar quality only, so its latency does not place it on our trade-off.}
\label{tab:e7}
\end{table}


\section{Overlap Study: Tests}
\label{app:v8}

The overlap study's analysis plan (internal) was fixed before any of its test decodes; it fixes the systems, the overlap
metrics and their thresholds, the causal relabelling rule and twelve confirmatory tests (Table~\ref{tab:reg8}). The
speaker-blind models are trained like \sys{} and \sys{}-Simul with the speaker tokens zeroed. Relabelling \sys{}'s own
output with the diarizer, keeping its text, separates labels from translation: on overlapped units its own labels score
\num{v8:ov:en2zh:O2:ov} against \num{v8:ov:en2zh:O2-relabel:ov} with diarizer labels (En$\rightarrow$Zh), while with
diarizer labels on both sides its text and the speaker-blind model's are close
(\num{v8:ov:en2zh:O2-relabel:ov} vs.\ \num{v8:ov:en2zh:O2-nosd-relabel:ov}).

\begin{table*}[t]
\centering\small
\setlength{\tabcolsep}{4pt}
\begin{tabular}{lrrrrrrrrr}
\toprule
 & \multicolumn{3}{c}{En$\rightarrow$Zh} & \multicolumn{3}{c}{Zh$\rightarrow$En} & \multicolumn{3}{c}{AliMeeting} \\
\cmidrule(lr){2-4}\cmidrule(lr){5-7}\cmidrule(lr){8-10}
System & all & overlap & CR & all & overlap & CR & all & overlap & CR \\
\midrule
\multicolumn{10}{l}{\emph{Reference windows, offline}} \\
Cascade, same data & 36.13 & 26.21 & 89.3 & 21.96 & 14.36 & 87.4 & 15.70 & 9.79 & 93.3 \\
Speaker-blind + relabel & 41.39 & 30.98 & 51.8 & 24.36 & 16.90 & 50.6 & 16.27 & 12.04 & 58.0 \\
\sys{} (offline) & 42.06 & 34.56 & 76.3 & 25.43 & 19.72 & 79.0 & 15.19 & 12.10 & 79.2 \\
\multicolumn{10}{l}{\emph{Whole meetings, online, 1\,s reads}} \\
Speaker-blind Simul + causal relabel & 35.22 & 27.17 & 65.5 & 19.33 & 9.06 & 54.4 & 12.67 & 8.55 & 63.0 \\
\sys{}-Simul & 40.09 & 35.71 & 82.3 & 24.53 & 17.73 & 80.2 & 11.47 & 9.18 & 80.2 \\
\bottomrule
\end{tabular}
\caption{Overlap-aware evaluation (\S\ref{sec:eval-overlap}): tcpBLEU over all units and over overlapped units, and
concurrent-speaker recall (CR, \%). Relabelling attaches one diarizer speaker to each translated segment. AliMeeting references: machine translations of the human transcripts (Appendix~\ref{app:data}).}
\label{tab:ovmain8}
\end{table*}

\input{generated/clean/tab_v8_reg}

\section{Analysis Plans}
\label{app:prereg}

Each study followed an internal analysis plan written before its test decoding and scoring, with every change recorded
as a dated amendment before the results it touches. The plans are timestamped only in a private repository, so
reviewers cannot verify them; we call their tests \emph{confirmatory} but do not rely on them as public
pre-registrations. The plans fix the meeting lists (AMI test EN2002a--d, ES2004a--d, IS1009a--d, TS3003a--d; all 20
AliMeeting test meetings; the eight AliMeeting evaluation meetings), the systems, the confirmatory contrasts with their
Holm families and the development rules for choosing checkpoints.
